\documentclass[11pt]{article}

\usepackage[margin=1.1in]{geometry}
\usepackage{amsmath}
\usepackage{amssymb}
\usepackage{amsthm}
\usepackage{booktabs}
\usepackage{microtype}
\usepackage[round,authoryear]{natbib}
\usepackage[colorlinks=true,allcolors=blue]{hyperref}

\newtheorem{proposition}{Proposition}

\title{\textbf{Exact Posterior Probability of Optimality\\ for
Factor-Gaussian Beliefs}\\[0.4em]
\large The Integral That Thompson-Style Methods Approximate,\\
Computed in Closed Form at Scale}
\author{Peter Cotton}
\date{\today}

\begin{document}
\maketitle

\begin{abstract}
\noindent The posterior probability that each alternative is best
drives stopping rules, budget allocation and Thompson sampling, and is
widely treated as intractable. Under common random numbers, and under
per-arm Gaussian observations with a factor prior, the posterior
covariance is exactly a low-rank factor term plus a diagonal at every
step. For such a covariance the whole probability vector is computable
exactly, for all $n$ alternatives at once, in time linear in $n$. In a
ranking-and-selection experiment, stopping rules built on marginal
means and variances spend up to three and a half times as many replicates
as the exact rule to reach the same error. In a correlated bandit, sampling
from the exact vector matches Thompson sampling, while a variational
policy built on marginals pays more regret.
\end{abstract}

\section{Introduction}

Let $\theta \in \mathbb{R}^n$ be the unknown values of $n$
alternatives, such as simulated systems, bandit arms or candidate
molecules, and let $\mathcal{D}$ be the data. The posterior
probability of optimality is the vector
\begin{equation}
p_i \;=\; \Pr\!\left(\theta_i = \max_j \theta_j \,\middle|\,
\mathcal{D}\right), \qquad i = 1, \dots, n.
\label{eq:pom}
\end{equation}

Sequential decision rules use this vector directly. A stopping rule
halts when $\max_i p_i$ reaches a confidence target. A budget
allocator spends where several $p_i$ are comparable. Thompson sampling
plays an action with probability $p_i$ without computing it. A policy
with constraints, or one that must be derandomized, needs the vector
itself, because one posterior sample does not determine it.

The Bayesian reinforcement learning literature describes
Equation~\eqref{eq:pom} as intractable. \citet{tarbouriech2023} call
it the key object of the field and write that computing it ``involves
computing several complicated integrals with respect to the posterior
and is intractable in most cases.'' Their method, VAPOR, replaces it
with a variational program over occupancy measures.

\citet{menet2026} fine-tune a large language model toward the same
object under a Gaussian-process posterior. The tractable policy they
target uses only the marginal means and marginal standard deviations
of that posterior. They note that the guarantee available for this
family assumes ``the worst-case bandit with independent arms.''

In ranking and selection, procedures bound the probability of correct
selection with Bonferroni-type inequalities or estimate it by
posterior sampling. The joint vector is not evaluated.

In several of these settings the posterior covariance has the exact
form
\begin{equation}
\Sigma \;=\; W W' + \mathrm{diag}(d), \qquad W \in
\mathbb{R}^{n \times \rho}, \quad \rho \ll n.
\label{eq:factor}
\end{equation}
This form is not a modelling approximation. It follows from how the
data are generated. For a covariance of this form,
Equation~\eqref{eq:pom} is computable exactly for every $i$ at once,
in time linear in $n$.

The computation is the pricing pass for correlated contests in
\citet{cotton2026}, which also develops its derivatives and its
inverse and benchmarks it to $n = 10^6$. This paper contributes three
things. Section~\ref{sec:posteriors} shows that two standard
posteriors keep the form of Equation~\eqref{eq:factor} after every
update. Section~\ref{sec:surrogates} identifies the approximations in
current use and the identity that ties the exact vector to Thompson
sampling. Sections~\ref{sec:stopping} and~\ref{sec:bandit} measure
what the exact vector changes in a stopping problem and in a bandit.

\section{Computing the vector for a factor-form covariance}
\label{sec:computation}

Write $\theta = m + W z + \varepsilon$, with $z \sim N(0, I_\rho)$ and
independent $\varepsilon_i \sim N(0, d_i)$. Given $z$, the
alternatives are independent, with means $m_i + w_i' z$ and variances
$d_i$. Let $F_{i|z}$ and $f_{i|z}$ be the conditional distribution and
density functions. Then
\begin{equation}
p_i \;=\; \int_{\mathbb{R}^\rho} \phi_\rho(z) \int_{\mathbb{R}}
f_{i|z}(x) \prod_{j \neq i} \bigl(1 - F_{j|z}(x)\bigr)\, dx \, dz .
\label{eq:integral}
\end{equation}

The inner integrands for different $i$ share the product
$G_z(x) = \prod_j (1 - F_{j|z}(x))$. Dividing out the $i$th term gives
every $p_i$ from the same product:
\begin{equation}
p_i \;=\; \mathbb{E}_z \int f_{i|z}(x)\,
\frac{G_z(x)}{1 - F_{i|z}(x)} \, dx .
\label{eq:cavity}
\end{equation}
On a lattice of $L$ points in $x$, with $Q$ quadrature nodes in $z$,
the whole vector costs $O(nLQ)$.

The numerical details are in \citet{cotton2026}: lattice placement,
products taken in logs, and the tails. That implementation is used
here without change. Every exact vector reported below was checked
against argmax frequencies from $2 \times 10^5$ Monte Carlo draws. The
largest total-variation distance in those checks is $0.0042$, which is
at the Monte Carlo noise level for that sample size.

\section{Two posteriors with exact factor form}
\label{sec:posteriors}

\subsection{Common random numbers}
\label{sec:crn}

Ranking and selection compares $n$ simulated systems. Under common
random numbers, replicate $r$ runs every system on the same random
scenario:
\begin{equation}
Y_{ir} \;=\; \theta_i + v_i' F_r + \varepsilon_{ir}, \qquad
F_r \sim N(0, I_\rho), \quad \varepsilon_{ir} \sim N(0, \sigma_i^2).
\end{equation}
The scenario effect $F_r$ is the same for every system. The loadings
$v_i$ and noise scales $\sigma_i$ are known properties of the
simulation. The unknown is $\theta$, with prior
$\theta \sim N(0, s_0^2 I)$.

The covariance of one replicate across systems is
\begin{equation}
\Lambda \;=\; V V' + D, \qquad D = \mathrm{diag}(\sigma_1^2, \dots,
\sigma_n^2).
\end{equation}
After $n_r$ replicates with mean vector $\bar{Y}$, the posterior of
$\theta$ is Gaussian with precision
\begin{equation}
A \;=\; s_0^{-2} I + n_r \Lambda^{-1}.
\label{eq:precision}
\end{equation}

The Woodbury identity gives
\begin{align}
\Lambda^{-1} &= D^{-1} - D^{-1} V M^{-1} V' D^{-1},
& M &= I_\rho + V' D^{-1} V.
\label{eq:wood1}
\end{align}
With $U = D^{-1} V$ and $S = n_r M^{-1}$, the precision becomes
\begin{align}
A &= \mathrm{diag}(a) - U S U', & a_i &= s_0^{-2} + n_r \sigma_i^{-2}.
\label{eq:diagminus}
\end{align}
This is a positive diagonal minus a positive semidefinite term of
rank $\rho$.

Applying the Woodbury identity again, to
Equation~\eqref{eq:diagminus}, gives
\begin{align}
A^{-1} &= \mathrm{diag}(a)^{-1} + \mathrm{diag}(a)^{-1} U \, C \, U'
\mathrm{diag}(a)^{-1},
& C &= \bigl(S^{-1} - U' \mathrm{diag}(a)^{-1} U\bigr)^{-1}.
\label{eq:wood2}
\end{align}

The matrix $C$ is positive definite. To see this, let
$K = \bigl(\begin{smallmatrix} \mathrm{diag}(a) & U \\ U' & S^{-1}
\end{smallmatrix}\bigr)$. The Schur complement of $S^{-1}$ in $K$ is
$\mathrm{diag}(a) - USU' = A$, which is positive definite, and
$S^{-1}$ is positive definite, so $K$ is positive definite. The Schur
complement of $\mathrm{diag}(a)$ in $K$ is $C^{-1}$, so $C^{-1}$ is
positive definite too.

The posterior covariance is therefore
\begin{equation}
\Sigma_{\text{post}} \;=\; \mathrm{diag}(d) + W W', \qquad
d_i = a_i^{-1}, \quad
W = \mathrm{diag}(a)^{-1} U \, C^{1/2},
\label{eq:postform}
\end{equation}
exactly, for every $n_r$. The rank of $W$ is $\rho$, the number of
shared scenario streams, not the number of systems. The posterior mean
is $m = A^{-1} h$ with $h = n_r \Lambda^{-1} \bar{Y}$, and both
factors are available from Equations~\eqref{eq:wood1}
and~\eqref{eq:wood2}. No step approximates. The posterior has factor
form because the scenario noise has factor form.

\subsection{Per-arm observations with a factor prior}
\label{sec:bandit-posterior}

In a finite Gaussian bandit the correlation sits in the prior on the
arm means: $\theta \sim N(0, \Sigma_0)$ with $\Sigma_0 = V V' +
\mathrm{diag}(d_0)$. Linear models $\theta_i = v_i' \beta + e_i$
produce this form, and so do inducing-point Gaussian-process
posteriors and the heteroscedastic output layers of large classifiers.
Pulling arm $a$ returns $\theta_a + N(0, \sigma_n^2)$.

With per-arm pull counts $c$, the posterior precision is
\begin{equation}
A \;=\; \Sigma_0^{-1} + \sigma_n^{-2}\,\mathrm{diag}(c).
\end{equation}
Equation~\eqref{eq:wood1} applied to $\Sigma_0$ puts $A$ in the form
of Equation~\eqref{eq:diagminus}, with $a_i = d_{0i}^{-1} + c_i
\sigma_n^{-2}$, $S = M_0^{-1}$ and $M_0 = I_\rho + V'
\mathrm{diag}(d_0)^{-1} V$. The rest of the argument is unchanged. The
posterior is exactly factor plus diagonal after every pull, with the
same rank $\rho$ as the prior.

\section{Approximations in current use, and Thompson sampling}
\label{sec:surrogates}

Two computable rules currently stand in for Equation~\eqref{eq:pom}.
Both use only the marginal means $m_i$ and marginal standard
deviations $s_i = \sqrt{d_i + \|w_i\|^2}$.

The \emph{independence approximation} treats the posterior as
$\prod_i N(m_i, s_i^2)$. \citet{menet2025} compute its probabilities
of being best with a threshold $\kappa^*$ chosen so that
$\sum_i \Phi((m_i - \kappa^*)/s_i) = 1$.

The \emph{VBOS policy} is the variational Bayesian optimistic sampling
rule of \citet{odonoghue2021}, extended by \citet{tarbouriech2023}. In
the near-closed form of \citet{menet2026} it is
\begin{equation}
\tilde{\pi}_i \;=\; v\!\left(\frac{m_i - \kappa^*}{s_i}\right),
\qquad v(c) = \exp\!\Bigl(-\tfrac{1}{8}\bigl(\sqrt{c^2 + 4} -
c\bigr)^2\Bigr),
\label{eq:vbos}
\end{equation}
with $\kappa^*$ again chosen so the vector sums to one.

The VBOS policy departs from $p$ on purpose, and its regret analysis
depends on that departure. It is a policy, not an estimate of $p$.
Its distance from $p$ is reported below because the fine-tuning
literature uses it as the trainable substitute for $p$.

Thompson sampling draws $\tilde{\theta}$ from the posterior and plays
the argmax of the draw. Its probability of playing arm $i$ at any
history is therefore $p_i$ at that history. Lemma~8 of
\citet{tarbouriech2023} states the general version: the expected
occupancy of Thompson sampling equals the posterior probability of
optimality.

Sampling an arm from the exact vector, called \emph{exact matching}
below, is therefore Thompson sampling in distribution. The two must
have the same expected regret, and Section~\ref{sec:bandit} checks
that they do. The exact vector adds a computed object, which a
stopping rule, an allocator or a derandomized policy can use where a
single draw cannot.

\section{Stopping under common random numbers}
\label{sec:stopping}

The experiment uses the posterior of Section~\ref{sec:crn} with
$n = 100$ systems, $\rho = 2$ scenario streams and prior scale
$s_0 = 1$. True means are drawn from the prior. Every system has
replicate variance one, so each replicate carries the same marginal
information in every regime. The \emph{factor share} is the fraction
of that variance that comes from the shared scenario.

The regimes differ only in how systems load on the scenario streams:
\begin{itemize}
\item \emph{Independent}: factor share zero. This is the control.
\item \emph{Aligned}: every system loads on the same stream, with
factor share $0.6$. Differences between systems then vary less than
their marginal variances suggest.
\item \emph{Opposed}: half the systems load positively and half
negatively on one stream, with factor share $0.6$. Differences between
systems in opposite halves then vary more than their marginal
variances suggest.
\end{itemize}

Three rules compete. Each computes its own estimate of $p$ after
every replicate, stops the first time its largest entry reaches
$1 - \delta$ with $\delta = 0.05$, and selects that entry. The exact
rule uses the computation of Section~\ref{sec:computation}. The other
two use the independence approximation and the VBOS policy.

\begin{table}[h]
\centering
\begin{tabular}{llccc}
\toprule
regime & rule & mean replicates & realized PCS & censored of 200 \\
\midrule
independent & exact & 141.2 & 0.970 & 35 \\
            & independence & 207.0 & 0.985 & 65 \\
            & VBOS  & 271.5 & 1.000 & 102 \\
\addlinespace
aligned     & exact & \phantom{0}69.7 & 0.958 & \phantom{0}9 \\
            & independence & 175.6 & 1.000 & 60 \\
            & VBOS  & 240.1 & 1.000 & 91 \\
\addlinespace
opposed     & exact & 106.3 & 0.944 & 21 \\
            & independence & 171.7 & 0.979 & 57 \\
            & VBOS  & 239.0 & 0.991 & 88 \\
\bottomrule
\end{tabular}
\caption{Stopping at nominal $95\%$ confidence. Each regime has 200
replications and a budget of 400 replicates. A run that reaches the
budget without stopping is censored and counts 400 toward the mean.
Realized PCS is the fraction of stopped runs that selected the true
best system.}
\label{tab:stopping}
\end{table}

Table~\ref{tab:stopping} shows that every rule held its nominal error.
The exact rule's $0.944$ in the opposed regime is within binomial noise
of $0.95$, with 179 stopped runs and a standard error of $0.016$.

The two marginal-only rules held the error by spending more. They used
$1.5$ to $3.4$ times the replicates of the exact rule and hit the
budget far more often. The gap is largest in the aligned regime. There
the leading systems' differences are tightly coupled, and the exact
rule stops after $69.7$ replicates on average, against $175.6$ and
$240.1$.

The marginal-only estimates are far from the exact vector along the
way. Before stopping in the aligned regime, the independence
approximation is at median total-variation distance $0.123$ from the
exact vector, with ninetieth percentile $0.174$. The VBOS policy is at
$0.141$, with ninetieth percentile $0.206$.

The distances do not vanish in the independent regime, where they are
$0.063$ and $0.096$. The independence approximation uses a shared
threshold and is biased even when the posterior is independent. The
VBOS policy departs from $p$ by design. The exact rule's estimate
matched Monte Carlo to within total-variation distance $0.0015$
throughout.

\section{A correlated Gaussian bandit}
\label{sec:bandit}

The bandit uses the posterior of Section~\ref{sec:bandit-posterior}
with $n = 50$ arms, prior variance one per arm, observation noise
$\sigma_n = 1$, horizon 200 and 100 replications. The prior has rank
two. In the \emph{independent} regime it has no factor term. In the
\emph{clusters} regime each half of the arms loads on its own factor,
which carries $0.75$ of the prior variance.

Four policies act on the same posterior: Thompson sampling, exact
matching, sampling from the VBOS policy, and sampling from the
independence approximation. Table~\ref{tab:bandit} reports cumulative
Bayes regret and, at the horizon, whether the arm with the highest
posterior mean is the true best and the shortfall of that arm's value
below the best. Its last column is the total-variation distance
between each policy's action distribution and the exact vector at
mid-horizon, along that policy's own trajectory.

\begin{table}[h]
\centering
\begin{tabular}{llcccc}
\toprule
regime & policy & regret & identification & shortfall &
TV to exact $p$ \\
\midrule
independent & Thompson & 172.1 & 0.72 & 0.073 & 0 \\
            & exact matching & 173.9 & 0.75 & 0.055 & 0 \\
            & VBOS & 190.3 & 0.76 & 0.061 & 0.159 \\
            & independence & 179.2 & 0.79 & 0.052 & 0.057 \\
\addlinespace
clusters    & Thompson & 122.7 & 0.61 & 0.097 & 0 \\
            & exact matching & 124.8 & 0.51 & 0.113 & 0 \\
            & VBOS & 135.9 & 0.59 & 0.105 & 0.129 \\
            & independence & 128.7 & 0.51 & 0.121 & 0.061 \\
\bottomrule
\end{tabular}
\caption{Mean over 100 replications of cumulative Bayes regret, the
rate at which the highest posterior mean identifies the true best arm
at the horizon, and the shortfall of that arm. The last column is the
total-variation distance from the exact vector at mid-horizon. It is
zero for Thompson sampling by Lemma~8 and for exact matching by
construction.}
\label{tab:bandit}
\end{table}

Thompson sampling and exact matching differ by one to two percent of
cumulative regret in both regimes, within replication noise, as
Lemma~8 requires. The VBOS policy pays ten to eleven percent more
regret than Thompson sampling, and its action distribution is at
total-variation distance $0.13$ to $0.16$ from the exact vector. The
independence approximation pays four to five percent more.

The identification rates do not separate the policies. Their
differences are within two binomial standard errors at 100
replications. The regret and distance columns carry the comparison.

\section{Scope}
\label{sec:limits}

The derivations of Section~\ref{sec:posteriors} are exact because the
observations there are generated with factor structure. When scenario
effects enter a simulation nonlinearly, a low-rank model of $\Lambda$
is an approximation, and the posterior inherits its error. When a
Gaussian-process kernel is dense with a short length scale, a factor
fit has to precede the computation, and its residual is the cost.

Under simple-regret loss, the one-step value of a measurement is the
knowledge-gradient quantity \citep{frazier2008}. \citet{tuisov2026}
shows that this is the normative myopic value of computation, and
notes that the exact dynamic version ``is itself generally
intractable.'' The method here does not compute that dynamic value.
For factor-Gaussian beliefs it makes the myopic quantities and the
stopping certificates exact.

The regret analysis of \citet{odonoghue2021} depends on the VBOS
tilt. Its measured distance from $p$ quantifies the approximation, but
does not by itself make a case against the tilt.

Thompson sampling needs one posterior sample and no vector. The exact
vector matters where a sample is not enough: stopping, allocation,
constraints and derandomization. Both experiments above are of that
kind.

\begin{thebibliography}{9}

\bibitem[Cotton(2026)]{cotton2026} P.~Cotton. Scalable inversion of
contests with correlated performances, including softmax and
multinomial probit. arXiv:2609.01133, 2026. Also SSRN 7307363.

\bibitem[Frazier et~al.(2008)Frazier, Powell, and Dayanik]{frazier2008}
P.~I. Frazier, W.~B. Powell, and S.~Dayanik. A knowledge-gradient
policy for sequential information collection. \emph{SIAM Journal on
Control and Optimization}, 47(5):2410--2439, 2008.

\bibitem[Menet et~al.(2025)Menet, Huebotter, Kassraie, and
Krause]{menet2025} N.~Menet, J.~Huebotter, P.~Kassraie, and
A.~Krause. Efficiently estimating Gaussian probability of maximality.
In \emph{AISTATS}, 2025. arXiv:2501.13535.

\bibitem[Menet et~al.(2026)Menet, Terzi\'c, Hersche, Krause, and
Rahimi]{menet2026} N.~Menet, A.~Terzi\'c, M.~Hersche, A.~Krause, and
A.~Rahimi. Thompson sampling via fine-tuning of LLMs. In
\emph{ICLR}, 2026. arXiv:2510.13328.

\bibitem[O'Donoghue and Lattimore(2021)]{odonoghue2021}
B.~O'Donoghue and T.~Lattimore. Variational Bayesian optimistic
sampling. In \emph{Advances in Neural Information Processing Systems
34}, 2021. arXiv:2110.15688.

\bibitem[Tarbouriech et~al.(2023)Tarbouriech, Lattimore, and
O'Donoghue]{tarbouriech2023} J.~Tarbouriech, T.~Lattimore, and
B.~O'Donoghue. Probabilistic inference in reinforcement learning done
right. In \emph{Advances in Neural Information Processing Systems
36}, 2023. arXiv:2311.13294.

\bibitem[Tuisov(2026)]{tuisov2026} A.~Tuisov. Search as computation
allocation. arXiv:2607.27871, 2026.

\end{thebibliography}

\end{document}
